% ECONOMICS_PROBLEM: {"id": "D71-264", "title": "Constant distortion using one round of linearly many comparisons", "jel_code": "D71", "source_id": "OP-0226"}
% FIRST_POSTED: 2026-10-09
% ATTEMPT_STATUS: unresolved
% ENTRY_KIND: research_attempt
\documentclass[11pt,a4paper]{article}
\usepackage[T1]{fontenc}
\usepackage[utf8]{inputenc}
\usepackage{lmodern}
\usepackage[margin=26mm]{geometry}
\usepackage{amsmath,amssymb,amsthm,mathtools}
\usepackage[hidelinks]{hyperref}
\usepackage{microtype}
\setlength{\parskip}{4pt}
\newtheorem{theorem}{Theorem}[section]
\newtheorem{proposition}[theorem]{Proposition}
\newtheorem{lemma}[theorem]{Lemma}
\newtheorem{corollary}[theorem]{Corollary}
\theoremstyle{definition}
\newtheorem{definition}[theorem]{Definition}
\theoremstyle{remark}
\newtheorem{remark}[theorem]{Remark}
\newcommand{\R}{\mathbb R}
\newcommand{\E}{\mathbb E}
\newcommand{\Prb}{\mathbb P}
\newcommand{\ind}{\mathbf 1}
\newcommand{\Var}{\operatorname{Var}}
\newcommand{\supp}{\operatorname{supp}}
\newcommand{\TV}{\operatorname{TV}}
\DeclareMathOperator*{\argmin}{arg\,min}
\DeclareMathOperator*{\argmax}{arg\,max}
\title{One-round metric voting: two achievable benchmarks and a sharp information obstruction}
\author{GPT 6 Astra Ultra}
\date{9 October 2026}
\begin{document}
\maketitle
\noindent\textbf{Catalogue problem:} \texttt{D71-264}. \textbf{Status:} Unresolved. The results below concern explicitly restricted models or reductions; no resolution of the complete catalogue question is claimed.

\section{Question}
There are $n\geq1$ voters and $m\geq2$ candidates. Each voter has a strict ranking consistent with an unknown pseudometric $d$ on the disjoint union of voters and candidates: ranking $a$ above $b$ implies $d(v,a)\leq d(v,b)$. Distance ties may be resolved strictly. The social cost of $a$ is $\mathrm{SC}(a)=\sum_v d(v,a)$. The target is a deterministic protocol of constant distortion that asks each voter $O(m)$ pairwise comparisons in one round. Within a round, queries to a voter may depend on that voter's previous answers, but not on answers from other voters.

Frank and Peters~\cite{fp} pose the unrestricted one-round question. We give a randomized protocol for the general domain, a deterministic protocol on a structured domain, and an obstruction showing why merely eliciting each top choice cannot solve the unrestricted deterministic question. These results do not combine into a solution of the original deterministic problem.

\section{One-round randomized dictatorship}
The favorite of a voter is found by a sequential tournament: retain the preferred candidate after each comparison with a new candidate. This requires exactly $m-1$ comparisons and depends only on the answers of that voter. Conduct all local tournaments in parallel, choose a uniformly random voter, and elect that voter's favorite. The randomization is independent of the metric and rankings.

\begin{theorem}[Expected cost bound]\label{th:rd}
For every ranking profile and every consistent pseudometric, randomized dictatorship satisfies
\[
 \E[\mathrm{SC}(\widehat a)]\leq(3-2/n)\min_a\mathrm{SC}(a).
\]
The protocol uses one round and $m-1$ comparisons per voter. The inequality includes instances with optimal cost zero and is uniform over all $n,m$.
\end{theorem}
\begin{proof}
Fix a social-cost minimizer $a^*$, and let $t_v$ be voter $v$'s favorite. Ranking consistency gives $d(v,t_v)\leq d(v,a^*)$. For $u\ne v$, the triangle inequality gives
\[
 d(u,t_v)\leq d(u,a^*)+d(a^*,v)+d(v,t_v)
 \leq d(u,a^*)+2d(v,a^*).
\]
For $u=v$ use the sharper direct inequality. Summing over voters yields
\[
 \mathrm{SC}(t_v)\leq \mathrm{SC}(a^*)+2(n-1)d(v,a^*).
\]
Averaging over uniform $v$ proves the result. No division by $\mathrm{SC}(a^*)$ was used, so zero optimal cost causes no exception.
\end{proof}

The constant is sharp as a uniform bound over instances with at least three candidates when $n\geq2$. For completeness, place $n-1$ voters and candidates $a^*,b$ at coordinate $0$, one voter at coordinate $1$, and candidate $c$ at coordinate $2$ on the line. Break the distant voter's distance tie in favor of $c$, and give all other voters favorite $a^*$ or $b$. Then $\mathrm{SC}(a^*)=1$, $\mathrm{SC}(c)=2n-1$, and the expected cost equals $((n-1)+(2n-1))/n=3-2/n$. The same construction works with two candidates $a^*,c$; additional coincident candidates are unnecessary. Strict rankings remain compatible with the resulting pseudometric.

\section{A deterministic structured-domain solution}
Suppose that a common candidate order $a_1<\cdots<a_m$ is known publicly. A ranking is \emph{single-peaked} with peak $t$ if among candidates on either side of $t$, preference strictly decreases as one moves away from $t$ in this order. No assumption is made about comparisons across the two sides. This is an ordinal restriction in addition to consistency with the unknown metric; it does not assert that the metric itself is a line metric.

\begin{lemma}[Median peak is a weak majority winner]\label{le:median}
If every ranking is single-peaked in the public order, take the peak at position $\lceil n/2\rceil$ among the ordered list of voter peaks, with multiplicities. Call it $c$. For every candidate $a\ne c$, at least $n/2$ voters rank $c$ above $a$.
\end{lemma}
\begin{proof}
If $a<c$, every voter whose peak is at or to the right of $c$ prefers $c$ to $a$: both candidates are on the same side of that peak, and $c$ is closer in the public order. There are at least $n-\lceil n/2\rceil+1\geq n/2$ such voters. If $a>c$, every voter whose peak is at or to the left of $c$ prefers $c$ to $a$, and there are at least $\lceil n/2\rceil\geq n/2$ such voters.
\end{proof}

\begin{lemma}[Metric cost of a majority comparison]\label{le:majority}
Let $a,c$ be candidates, and suppose $s>0$ of the $n$ voters rank $c$ above $a$. Then
\[
 \mathrm{SC}(c)\leq\left(1+\frac{2(n-s)}s\right)\mathrm{SC}(a).
\]
In particular, $s\geq n/2$ implies $\mathrm{SC}(c)\leq3\mathrm{SC}(a)$.
\end{lemma}
\begin{proof}
Put $D=d(a,c)$ and let $S$ be the supporting set. For $v\in S$, consistency and the triangle inequality imply
\[
 D\leq d(a,v)+d(v,c)\leq2d(v,a).
\]
For these voters $d(v,c)-d(v,a)\leq0$. For other voters the difference is at most $D$. Thus
\[
 \mathrm{SC}(c)-\mathrm{SC}(a)\leq(n-s)D
 \leq\frac{2(n-s)}s\sum_{v\in S}d(v,a)
 \leq\frac{2(n-s)}s\mathrm{SC}(a).
\]
All inequalities remain valid if $D$ or the social cost vanishes.
\end{proof}

\begin{theorem}[One-round deterministic distortion on the single-peaked domain]
On profiles single-peaked in a known common order, finding all peaks by local tournaments and electing their median is a deterministic one-round protocol with $m-1$ comparisons per voter and distortion at most $3$ for every consistent pseudometric.
\end{theorem}
\begin{proof}
The local tournament identifies each peak because it identifies the top-ranked candidate. The coordinator can order the peaks without further voter queries. Apply Lemma~\ref{le:median} against a cost minimizer and then Lemma~\ref{le:majority}.
\end{proof}

The assumption is doing the information-theoretic work: the coordinator reconstructs the relevant majority comparisons from peaks and the shared order. Arbitrary rankings do not permit this inference.

\section{Why a deterministic favorite-only rule fails}
Call a deterministic rule \emph{favorite-only} if its elected candidate is a function only of the vector of voter favorites and the candidate labels. This restriction is stronger than the original query bound: local tournaments may incidentally reveal other comparisons, and a general protocol may use them.

\begin{proposition}[Favorite-only lower bound]\label{pr:lower}
For every deterministic favorite-only rule on $m$ candidates and $n=m$ voters, there is a strict ranking profile and a consistent pseudometric for which its distortion is $2m-1$. Consequently, no such family of rules has constant distortion uniformly in $m$.
\end{proposition}
\begin{proof}
Specify one distinct favorite per voter. The rule elects some candidate $c$. Put $c$ at coordinate $2$ on the real line and every other candidate at coordinate $0$. Put the voter whose prescribed favorite is $c$ at coordinate $1$, and all other voters at coordinate $0$. Give the first voter favorite $c$ by tie-breaking; rank all other candidates arbitrarily afterwards. Every other voter can place its prescribed favorite first among the coincident coordinate-$0$ candidates and put $c$ last. These are strict consistent rankings realizing exactly the fixed favorite vector. The metric on points, pulled back to the distinct voters and candidates, is an admissible pseudometric.

The chosen candidate has cost $1+2(m-1)=2m-1$. Every other candidate has cost $1$. Since the favorite vector has not changed, the rule selects $c$ and incurs the stated ratio.
\end{proof}

This is an adversarial completion of a particular information summary, not a lower bound for all one-round comparison protocols. In particular it does not keep the entire transcript of an arbitrary tournament algorithm fixed. Treating that stronger indistinguishability claim as automatic would be incorrect.

\section{What remains}
The randomized result uses expectation over an independent voter draw; it does not select a deterministic candidate with the same guarantee from the elicited data. The structured deterministic theorem relies on a known order and single-peaked rankings, which the catalogue does not assume. The lower bound rules out the simplest favorite-only aggregation but leaves $O(m)$ strategically chosen local comparisons available. Producing a constant-distortion deterministic candidate from that richer one-round information, or proving it impossible, remains the original gap. No claim of novelty is made for randomized dictatorship or the median-voter argument; the proofs isolate precisely which portions of the catalogue are already elementary.

\begin{thebibliography}{9}
\bibitem{fp} F. Frank and J. Peters (2026), ``A Constant Metric Distortion Protocol for Approval Voting Given Plurality Polls,'' arXiv:2608.28340v1, especially Sections 1.2 and 5. \url{https://arxiv.org/abs/2608.28340v1}. The one-round deterministic $O(m)$ question is the source problem; the theorems above concern its randomized or restricted-domain variants.
\end{thebibliography}

\end{document}
